\section{Methods}

\subsection{Growth and city trajectories}
Fixed-member comparisons distinguish changes within observed cities from changes in archive coverage. Activity and facilities use the compatible-week measures in Equation~\eqref{eq:metrics}; their ratio evaluates activity relative to observed scale.

Growth comparisons use eight pre-defined common-city windows: the 2019--2023 cohort, four adjacent-year windows, and three continuous three-year windows.  A city enters a window only when the required endpoint and intermediate years have positive activity and facilities, and membership is then held fixed for the endpoint comparison.  For endpoints $0$ and $1$, define $x_i=\ln(N_{i1}/N_{i0})$, $y_i=\ln(A_{i1}/A_{i0})$, and $dI_i=y_i-x_i$.  The descriptive regimes are facility expansion with faster activity growth ($x_i>0$, $y_i>x_i$), expansion with slower activity growth ($x_i>0$, $0<y_i<x_i$), and expansion with activity contraction ($x_i>0$, $y_i<0$).  Other cases, including facility non-growth and exact ties, are retained separately.  A tolerance of $10^{-10}$ is used only to classify floating-point values at zero or on $y=x$.

Network activity and facility growth are calculated from sums of city-level $A$ and $N$, respectively, within each fixed membership; overall intensity is their ratio. Equal-city log changes and interquartile ranges describe the distribution across members. Facility-weighted overall intensity and equal-city distributions therefore answer different questions even within the same window.

To examine persistence, growth types are paired within cities across two windows using the same endpoint log changes and classification rules.  The primary comparison uses the non-overlapping 2020--2021 and 2022--2023 windows; the supplementary comparison uses 2021--2022 and 2022--2023, which share 2022.  Each city contributes its early and late regime labels only when both labels are defined under the pre-specified rules; unqualified or missing records are not assigned to ``Other.''  The transition matrix places the early label in columns and the late label in rows.  Its cell count is $n_{rc}$ and its displayed conditional percentage is $100n_{rc}/\sum_{r'}n_{r'c}$, so percentages sum down each early-category column; an empty column has no defined denominator.  The paired comparisons describe within-city changes in growth type. The non-overlapping comparison spans different pandemic-response periods; it does not attribute type changes to the pandemic, policy or network maturity.

\subsection{Recorded activity states and aggregate change}
To include cities with source-coded-zero activity, we form the four adjacent endpoint windows 2019--2020, 2020--2021, 2021--2022 and 2022--2023 independently.  A city enters a window when both endpoint records have finite $Q\geq0$ and finite positive $E$, $W$ and $N$, with $A=Q/W$ and $N=E/W$ defined at both endpoints.  The endpoint member set is fixed within that window; no weather-completeness, model-sample, separation or singleton restriction is imposed.  We classify the signs of endpoint activity as positive or source-coded zero: $++$ has $Q_0>0,Q_1>0$, $0+$ has $Q_0=0,Q_1>0$, $+0$ has $Q_0>0,Q_1=0$, and $00$ has $Q_0=Q_1=0$.  These groups identify changes in the sign of recorded activity; they carry no assumption about commissioning or shutdown.

For each state window, let $S_{A,t}=\sum_{i\in C}A_{it}$ and $S_{N,t}=\sum_{i\in C}N_{it}$, with aggregate recorded intensity $I^{\mathrm{agg}}_t=S_{A,t}/S_{N,t}$.  We report endpoint levels and changes for the full member set and, separately, for its $++$ subset.  The facility share of state group $g$ at endpoint $t$ is $s_{g,t}=\sum_{i\in g}N_{it}/S_{N,t}$; these shares sum to one over the four states at each endpoint.  Activity change is decomposed exactly as
\begin{equation}
\Delta S_A=\bigl(S_{A,++,1}-S_{A,++,0}\bigr)+S_{A,0+,1}-S_{A,+0,0}+0_{00}.
\label{eq:record-state-decomposition}
\end{equation}
The displayed state contributions are divided by the full-set baseline $S_{A,0}$ and expressed in percentage points, so they sum to the full-set percentage change when the baseline is nonzero.  A zero baseline leaves the relative change undefined rather than receiving an arbitrary offset.  The full member set and its positive subset each use their own endpoint totals. SI Section~\siref{sec:si-descriptive} gives complete state-specific levels and shares.



\subsection{Intensity contributions and recorded-week composition}
Changes in network intensity can arise within cities or from changes in their facility weights. To distinguish these sources, let $w_{it}=N_{it}/\sum_jN_{jt}$ and $I_t^{\mathrm{agg}}=\sum_iw_{it}I_{it}=\sum_iA_{it}/\sum_iN_{it}$.  The symmetric accounting decomposition is
\begin{equation}
 \Delta I^{\mathrm{agg}}=
 \sum_i\frac{w_{i0}+w_{i1}}{2}(I_{i1}-I_{i0})+
 \sum_i\frac{I_{i0}+I_{i1}}{2}(w_{i1}-w_{i0}).
 \label{eq:decomposition}
\end{equation}
The terms are, respectively, the within-city intensity component at average endpoint weight and the composition component at average endpoint intensity.  This is an accounting identity, not a behavioral or policy mechanism \citep{kitagawa1955}.  We retain each city's exact contribution,
\begin{equation}
 \begin{split}
 \mathrm{within}_i&=\frac{w_{i0}+w_{i1}}{2}(I_{i1}-I_{i0}),\\
 \mathrm{composition}_i&=\frac{I_{i0}+I_{i1}}{2}(w_{i1}-w_{i0}),\\
 \mathrm{total}_i&=\mathrm{within}_i+\mathrm{composition}_i
 =w_{i1}I_{i1}-w_{i0}I_{i0}.
 \end{split}
 \label{eq:city-contribution}
\end{equation}
City contributions are normalized by their window's initial aggregate intensity, set to 100. A city's own intensity change and its weighted contribution are evaluated separately because a change in facility share can reverse their signs. The largest five positive and negative contributors are distinguished from the remaining cities.

A second intensity decomposition identifies how recording frequency and positive-week activity combine with facility scale. At an endpoint with nonnegative compatible-week activity and at least one positive week, let $p$ be the fraction with positive recorded activity and $q^{+}$ its positive-week mean.  Then
\begin{equation}
 A=pq^{+},\qquad I=\frac{pq^{+}}{N},\qquad
 \Delta\ln I=\Delta\ln p+\Delta\ln q^{+}-\Delta\ln N.
 \label{eq:weeks}
\end{equation}
The identity is evaluated for the fixed 36-city 2019--2023 and 93-city 2022--2023 memberships.  Missing activity remains outside it; equal-city means satisfy the identity, whereas component medians need not add.

For the two focal contribution windows, the fixed-member comparison uses their 52-city intersection and recomputes weights at each endpoint. For $m$ negative contributors, $s_i$ is each absolute share of the negative total; we report $\mathrm{HHI}=\sum_i s_i^2$, its normalized form, and largest-five and largest-fifth shares.

\subsection{Annual conditional relationships}
Conditional models distinguish a positive facility--activity association from a proportional relationship: the relevant hypotheses are zero and unit elasticity, respectively. Annual models compare city-year activity at given facilities, exposure and controls. The baseline log-OLS model is
\begin{equation}
 \ln A_{it}=\alpha_i+\lambda_t+\beta\ln N_{it}+\gamma^{\top}Z_{it}+\varepsilon_{it},
 \label{eq:ols}
\end{equation}
where the positive-$Q$ sample supplies the log outcome.  The corresponding level model is estimated by Poisson pseudo-maximum likelihood (PPML):
\begin{equation}
 \mathbb{E}[Q_{it}\mid N_{it},W_{it},Z_{it}]
 =W_{it}\exp\{\alpha_i+\lambda_t+\beta\ln N_{it}+\gamma^{\top}Z_{it}\}.
 \label{eq:ppml}
\end{equation}
The first equation models a conditional mean on the log weekly-activity scale; the second models a conditional arithmetic mean for cumulative activity at observed exposure.  Different conditional means can produce different coefficients, so comparisons concern each model's estimand rather than agreement on a single universal elasticity.

Both models include city and calendar-year effects.  $Z_{it}$ contains five city-year weather controls (heating and cooling degree-days, precipitation, wind and humidity).  The positive model has 445 rows in 104 cities; the zero-inclusive PPML has 523 rows in 112 cities (453 positive, 70 source-coded zero), requiring finite nonnegative $Q$, positive $E,W$ and complete controls.  Unknown activity is never recoded as zero, so the supports answer related questions rather than one zero-inclusion contrast. Sample exclusions and the full support accounting are given in SI Section~\siref{sec:si-annual}.

A restricted cubic spline adds two basis terms at the empirical 5th, 35th, 65th and 95th percentiles of positive-support $\ln N$ \citep{harrell2015}.  Its derivative is a local elasticity, not a threshold or optimum; basis details, pointwise intervals and the curve are in SI Section~\siref{sec:si-annual}.  The two added terms are tested jointly with an $F(2,G-1)$ no-curvature test, distinct from $\beta=0$ or $\beta=1$.

Two-sided $t(G-1)$ tests cover $\beta=0$ and $\beta=1$; 10\% facility contrasts use $100(1.1^{\widehat\beta}-1)$ with transformed endpoints.  In PPML, $\ln W$ is an exposure offset, not a frequency weight \citep{silva2006}; city-clustered CR1 inference uses full design rank and $t(G-1)$ references \citep{cameron2015}.

\subsection{Quarterly dynamic conditional relationship}
The quarterly model distinguishes the contemporaneous relationship from additional association over the next three quarters. This finite horizon complements the annual scale comparison.  The quarterly estimation sample requires finite current $Q\geq0$, $E,W>0$, positive $N$ at four consecutive calendar quarters, and city-matched lags.  Lagged $Q$ may be coded zero; no quarterly weather is available or imputed.  Support exclusions and separation diagnostics are in SI Section~\siref{sec:si-quarterly} \citep{correia2026separation}.

The quarterly PPML specification includes three lags ($L=3$):
\begin{equation}
 \mathbb{E}[Q_{it}\mid \{x_{i,t-j}\}_{j=0}^{3},W_{it}]
 =W_{it}\exp\{\alpha_i+\tau_t+\beta_0x_{it}+\beta_1x_{i,t-1}
 +\beta_2x_{i,t-2}+\beta_3x_{i,t-3}\},\qquad x_{it}=\ln N_{it},
 \label{eq:quarterly-d3}
\end{equation}
where $\alpha_i$ are city effects and $\tau_t$ are complete calendar-quarter effects.  Static, one-lag and three-lag specifications use the same estimable support, offset and city-clustered inference.  An equivalent level-and-difference parameterization, complete coefficient paths and the COVID comparison are in SI Section~\siref{sec:si-quarterly}.  The COVID pair restricts this support through 2022 and adds the current index divided by 10; it uses province clustering and excludes 2023 because complete response-index exposure is unavailable.

For a model with maximum lag $L$, the endpoint elasticity is $c_L=\sum_{j=0}^{L}\beta_j$: $c_0$, $c_1$ and $c_3$ for the static, one-lag and three-lag specifications, respectively (Table~\ref{tab:quarterly_dynamic}). Within each specification, a stipulated persistent facility difference gives a conditional weekly-mean elasticity at horizon $L$, holding covariate and fixed-effect paths.  It is not cumulative activity or a causal impulse response.  The unit-elasticity test $c_3=1$ evaluates proportionality at the finite-horizon endpoint, whereas $c_3-c_0=0$ tests additional follow-through. Both use the full clustered covariance, with pointwise 95\% intervals and the corresponding 10\% facility transformation. Contrasts therefore incorporate covariance between coefficients rather than subtracting individual interval endpoints.

\subsection{Policy content, annual associations and governance timing}
Policy evidence links evaluation objects to prescribed operating responsibilities. Clause coding establishes what is required; annual associations relate those requirements to city outcomes; quarterly timing comparisons examine changes around dated governance records. Locally applicable requirements are coded at clause level.  The policy corpus comprises 601 locally applicable policy documents covering 176 cities. Of these, 207 documents covering 88 cities provide 466 clauses for the analysis of operational requirements.  Four non-exclusive functions are coded: information/supervision, operation/maintenance, continuing operation/payment and assessable or subsidy-linked conditions.  Counts describe prescribed responsibilities, not compliance; the content figure and administrative sequence use the 88-city and separate 95-city scopes, respectively. Coding definitions and scope intersections are detailed in SI Section~\siref{sec:si-policy}.

For each $Y\in\{A,N,I\}$, the annual policy models relate log outcomes to city/year effects, prior-year-end information/supervision $R$, operation/maintenance $S$, the previous-year construction-support document count $B$, and annual temperature, precipitation, wind and humidity controls.  For $B$, a policy is assigned to the later of its effective and publication years when both are known, to its effective year when only that year is known, and to no known year when its effective year is missing; the regressor counts documents assigned to year $t-1$, rather than accumulated policies. A second specification adds the current response index divided by 10.  Both use 325 city-year rows in 94 cities with province-clustered inference.  $R$ and $S$ are observed exposure, not implementation intensity, and are not joined to the quarterly model as a mediation chain. The complete policy and weather specifications are in SI Section~\siref{sec:si-associations}.

The 95-city timing registry codes construction support $B$, information/supervision $R$ and operation/maintenance responsibility $S$.  Date intervals define clean pre-policy, transitional and post-policy periods; year-only dates remain intervals, unknown coverage is not untreated, and first facility observation is not construction or commissioning.

The governance extension fits city and calendar-quarter effects to eligible clean-pre or conditionally never-observed rows, excludes transition/unknown rows, and uses no quarterly weather.  Supported outcomes are predicted, averaged within city and then across cities \citep{borusyak2024}.  The primary target contains 11 cities and 36 city-quarters; the four-lead facility diagnostic ($p=0.0008415$) uses a different reference and is descriptive rather than a net governance effect.  Full support rules and window decompositions are in SI Section~\siref{sec:si-governance}.
